\chapter{Porphyrins}

\section*{What Is This Test?}
Porphyrin testing measures porphyrin precursors and porphyrin intermediates in urine, blood, or stool to diagnose and classify the porphyrias, a heterogeneous group of rare inherited disorders of heme biosynthesis. The pathway: glycine + succinyl-CoA $\rightarrow \delta$-aminolevulinic acid (ALA) $\rightarrow$ porphobilinogen (PBG) $\rightarrow$ porphyrin intermediates $\rightarrow$ heme. Enzyme defects at any step cause accumulation of proximal intermediates, producing neurovisceral attacks (acute porphyrias) or cutaneous photosensitivity (cutaneous porphyrias). Test selection depends on clinical presentation: acute attacks require urine PBG and ALA; cutaneous symptoms require plasma/urine/fecal porphyrin fractionation.

\section*{Alternative Names}
Porphyrins; Porphobilinogen (PBG); Delta-Aminolevulinic Acid (ALA); Plasma Porphyrins; Urine Porphyrins; Fecal Porphyrins; Fractionated Porphyrins; Free Erythrocyte Protoporphyrin

\section*{Principle}
Quantification by high-performance liquid chromatography (HPLC) with fluorescence or mass spectrometry (LC-MS/MS) has replaced older solvent extraction methods. Specimen choice depends on suspected type: urine (24-hour) for PBG, ALA, coproporphyrin, uroporphyrin; whole blood (EDTA) for zinc protoporphyrin and free erythrocyte protoporphyrin; feces for protoporphyrin and coproporphyrin fractionation; plasma for total plasma porphyrin (cutaneous symptoms). Acidified urine protects from photodegradation.

\section*{History}
Porphyrin chemistry began with Hoppe-Seyler's 1871 identification of porphyrin pigments. Acute intermittent porphyria was clinically described by Waldenström in 1937. Modern enzymatic and genetic characterization of the eight porphyrias occurred through the late 20th century. Porphobilinogen deaminase deficiency was identified as the cause of acute intermittent porphyria in the 1980s.

\section*{Why This Test Is Ordered}
\begin{itemize}
  \item Acute neurovisceral attack: severe abdominal pain, vomiting, autonomic instability, motor neuropathy, mental status changes; sometimes triggered by drugs, fasting, hormones
  \item Cutaneous photosensitivity: painful skin fragility, blisters on sun-exposed areas
  \item Evaluation of abdominal pain of unknown cause with neurologic or psychiatric features
  \item Family screening in known mutation carriers
  \item Lead poisoning workup (elevated urinary ALA and coproporphyrin)
  \item Pre-symptomatic testing of at-risk relatives in confirmed porphyria families
\end{itemize}

\section*{Is This a Primary or Secondary Test}
Secondary, ordered when clinical suspicion for porphyria is raised; not a routine screening test.

\section*{How This Test Is Performed}
A 24-hour urine collection in dark container with 5 g sodium carbonate preservative protects from light; whole blood in EDTA; random fecal sample in opaque container. Specimens must be protected from light during collection and transport. Samples for PBG/ALA must be refrigerated and assayed within 48 hours.

\section*{What Preparation Is Needed}
None specific. Document medications taken over the preceding month (porphyrinogenic drugs, hormones, anticonvulsants). Stop triggering drugs before sampling if possible. Avoid alcohol, fasting, and strenuous exercise.

\section*{Contraindications and Precautions}
Light exposure denatures porphyrins: samples must be protected. False elevations from lead exposure, iron deficiency, liver disease, alcoholism, certain drugs. Differentiation of porphyrias requires fractionated porphyrins, not just a total porphyrin measurement.

\section*{Risks}
Venipuncture only.

\section*{What Happens After the Test}
Elevated urine PBG/ALA confirms acute porphyria. Subsequent fractionation in urine, feces, blood, plasma distinguishes acute intermittent porphyria, hereditary coproporphyria, variegate porphyria, porphyria cutanea tarda, erythropoietic protoporphyria, etc. Genetic testing confirms the specific gene defect.

\section*{Understanding Results}

\subsection*{Below Normal}
Normal PBG/ALA in the presence of acute symptoms excludes acute neurovisceral porphyria. Normal urine porphyrins does not exclude cutaneous porphyrias requiring plasma or fecal testing.

\subsection*{Above Normal}
\begin{itemize}
  \item \textbf{Acute intermittent porphyria (AIP):} Markedly elevated urine PBG ($>$ 10 mg/24 h), ALA elevated; reduced in remission but rarely normal
  \item \textbf{Hereditary coproporphyria:} Elevated urine PBG and ALA, increased fecal coproporphyrin III
  \item \textbf{Variegate porphyria:} Elevated urine PBG and ALA; markedly increased fecal protoporphyrin and coproporphyrin; plasma fluorescence at 626 nm
  \item \textbf{Porphyria cutanea tarda (PCT):} Normal PBG; elevated urine uroporphyrin $>$ coproporphyrin (greater than 8:1)
  \item \textbf{Erythropoietic protoporphyria (EPP):} Elevated free erythrocyte protoporphyrin and plasma protoporphyrin; urine normal
  \item \textbf{Lead poisoning:} Elevated urinary ALA and coproporphyrin III with normal PBG
\end{itemize}

\section*{Normal Results}
\begin{itemize}
  \item \textbf{24-h urine PBG:} 0--2 mg/24 h
  \item \textbf{24-h urine ALA:} 0--7 mg/24 h
  \item \textbf{Total urine porphyrins:} 20--120 $\mu$g/24 h
  \item \textbf{Plasma porphyrin:} $<$ 1 $\mu$g/dL
  \item \textbf{Fecal protoporphyrin:} 10--80 $\mu$g/g dry weight
  \item \textbf{Free erythrocyte protoporphyrin / Zinc protoporphyrin (ZPP):} $<$ 35 $\mu$g/dL
  \item \textbf{Elevated PBG:} $>$ 6 mg/24 h strongly suggests AIP, HCP, or VP during attack
\end{itemize}

\section*{Follow-up Tests}
\begin{itemize}
  \item \textbf{Porphyrin fractionation (urine, feces, plasma):} Confirms subtype
  \item \textbf{Enzyme activity (e.g., PBG deaminase in RBCs):} Confirms AIP in some cases
  \item \textbf{Genetic testing (HMRS, CPOX, PPOX, UROD, FECH genes)}: Provides definitive subtype and enables cascade family testing
  \item \textbf{Serum iron, ferritin, hepatitis C and HIV antibodies, alcohol history}: Evaluate precipitating factors in PCT
  \item \textbf{Liver function tests and abdominal imaging}: Assess hepatotoxicity or hepatic tumors (HCC risk increased in PCT, AIP, VP)
  \item \textbf{Avoid porphyrinogenic drugs}: Barbiturates, sulfonamides, rifampin, oral contraceptives; consult safe drug lists at porphyriafoundation.org
\end{itemize}

\section*{References}
\begin{enumerate}
  \item Bissell DM, Anderson KE, Bonkovsky HL. Porphyria. N Engl J Med. 2017;377(9):862-872.
  \item Anderson KE, et al. Recommendations for the diagnosis and treatment of the acute porphyrias. Ann Intern Med. 2005;142(5):369-379.
  \item Puy H, Gouya L, Deybach JC. Porphyrias. Lancet. 2010;375(9718):924-937.
  \item American Porphyria Foundation. Drug safety database. www.porphyriafoundation.org.
\end{enumerate}

\clearpage